%% kapa: exam=1 round=22 year=2011 period=2 form=A subject=accounting subject_ko=회계학 seq=21 number=21
%% source: Q-Net 감정평가사 기출문제 (cst00309) artlSeq=5036790 file=extracted/1008/5036790/2011년 제22회 감정평가사 1차시험 기출문제/2교시.hwp
%% conversion: kordoc 4.24.51 -> tools/gen1.py (auto); review=proofread-gpt-astra; flags=none
\begin{question}{21}
리스 회계처리에 관한 설명으로 옳지 않은 것은?
\begin{choices}
  \item 리스약정일은 리스계약일과 리스의 주요사항에 대한 계약당사자들의 합의일 중 이른 날이다.
  \item 리스기간개시일은 리스이용자가 리스자산의 사용권을 행사할 수 있게 된 날로 리스자산의 최초인식일이 된다.
  \item 리스기간 중에 리스자산의 소유권이 리스이용자에게 이전되는 경우에는 금융리스로 분류한다.
  \item 리스의 분류는 리스기간개시일을 기준으로 결정한다.
  \item 리스기간은 리스이용자가 자산을 리스하기로 약정을 맺은 해지불능기간과 리스이용자가 리스를 연장할 수 있는 선택권을 가지고 있으며, 리스이용자가 그 선택권을 행사할 것이 리스약정일 현재 거의 확실한 경우 그 추가기간을 포함한다.
\end{choices}
\end{question}
\answer{4}
